\SourceLeaf{315}%
Sobald sich der Wert von \(T\) bei Verwendung fünfstelliger Logarithmen nicht mehr ändert, berechnet man \(\lambda\) nach der Formel
\[
\lambda=r\sin T\sec B
\]
und führt dann eine neue Rechnung mit sieben Dezimalstellen durch. Dabei zieht man die logarithmischen Korrekturen mithilfe der Formeln
\[
\begin{aligned}
\log b&=\log\beta+\mu\lambda\lambda+\tfrac12\mu\tau\tau\\
\log t&=\log\tau+\mu rr+\tfrac12\mu\tau\tau
\end{aligned}
\]
heran und setzt \(B=B^0-\frac12b\) und \(T=T^0-\frac12t\). Eine weitere Wiederholung ergibt in der Regel dieselben oder kaum merklich veränderte Ergebnisse. Erst dann fügt man auch die Berechnung von \(l\) nach der Formel
\[
\log l=\log\lambda-\tfrac12\mu rr+\tfrac12\mu\lambda\lambda
\]
hinzu. Um an einem Beispiel zu zeigen, wie schnell die Annäherung verläuft (was hauptsächlich davon abhängt, wie klein \(r\) ist), gebe ich hier die wichtigsten Schritte der Rechnung vom Dreieckspunkt Brocken zum Punkt Inselsberg an. Das ist die längste Dreiecksseite im hannoverschen Dreiecksnetz, viel länger als die Seiten, die sonst bei trigonometrischen Vermessungen vorkommen.

Bei der nach den Grundlagen der ersten Abhandlung ausgearbeiteten konformen Abbildung auf die Kugeloberfläche beträgt die Breite des Brockens \(B^0=51^0\,46'\,3''\,6345\), das Azimut der Seite Brocken-Inselsberg beträgt \(T^0=5^0\,42'\,21''\,7704\), und der Logarithmus dieser Seite ist in Toisen \(=4{,}7353929\), in Einheiten des Radius \(=8{,}22018543\) oder in Bogensekunden, wie in unseren Formeln vorausgesetzt, \(\log r=3{,}5346106\). Setzt man zunächst \(T=5^0\,42'\), so erhält man
\[
\begin{aligned}
\beta&=3408''\\
B&=51^0\,17'\,40''\\
\tau&=424''
\end{aligned}
\]
und folglich den besseren Näherungswert
\[
T=5^0\,38'\,50''
\]
